\section{Posterior computation over ordered partitions}
\label{sec:exact-inference}

Once the prior-adjusted block array
$\widetilde A^{(r)}_{ij}=A^{(r)}_{ij}\gamma_u(i,j)$ has been constructed, the observation family no longer enters the partition recursions. This section gives the two distinct inference routes used by \BayesBreak: sum-product for posterior quantities and max-sum for a joint maximizing partition.

\begin{figure}[!htbp]
  \centering
  \resizebox{0.91\textwidth}{!}{
\begin{tikzpicture}[node distance=8mm and 13mm]
\node[secondarybox,text width=30mm] (scores) {Integrated segment log scores\\$\ell_{ij}$};
\node[primarybox,above right=of scores,text width=28mm] (sum) {Sum-product\\$\log\sum\exp$};
\node[primarybox,below right=of scores,text width=28mm] (map) {Max-sum\\$\max$};
\node[neutralbox,right=of sum,text width=34mm] (posterior) {$p(k\mid y)$, boundary marginals, signal moments};
\node[neutralbox,right=of map,text width=34mm] (maxpath) {Joint MAP partition with backtracking};
\draw[arr] (scores)--(sum);
\draw[arr] (scores)--(map);
\draw[arr] (sum)--(posterior);
\draw[arr] (map)--(maxpath);
\end{tikzpicture}
}
  \caption{The same local scores support two algebraically related but inferentially different routes. Sum-product integrates over paths; max-sum selects one globally maximizing path. Coordinatewise modes of marginal boundary distributions need not form that path.}
  \label{fig:semiring-routes}
\end{figure}

\subsection{Forward and backward partition sums}
For $k\in\{0,\ldots,k_{\max}\}$ and $j\in\{0,\ldots,n\}$, let $L_{k,j}$ be the integrated score of all $k$-block partitions of the prefix $(0,j]$. Let $R_{k,i}$ be the corresponding score for all $k$-block partitions of the suffix $(i,n]$. Set
\begin{align*}
L_{0,0}&=1,\qquad L_{0,j}=0\quad(j>0),\\
R_{0,n}&=1,\qquad R_{0,i}=0\quad(i<n).
\end{align*}
The sum-product recursions are
\begin{align}
L_{k+1,j}&=\sum_{h=k}^{j-1}L_{k,h}\widetilde A^{(0)}_{h,j},
\label{eq:forward}\\
R_{k+1,i}&=\sum_{h=i+1}^{n-k}\widetilde A^{(0)}_{i,h}R_{k,h}.
\label{eq:backward}
\end{align}
In implementation, Equations~\eqref{eq:forward}--\eqref{eq:backward} are evaluated in log space using \texttt{logsumexp}.

\begin{theorem}[Partition-sum identities and boundary marginals]
\label{thm:sum-product}
Under Assumptions~\ref{ass:offline} and~\ref{ass:factorized-prior}, for every admissible $k$,
\begin{equation}
L_{k,n}=R_{k,0}
=\sum_{0=t_0<\cdots<t_k=n}\prod_{q=1}^{k}\widetilde A^{(0)}_{t_{q-1},t_q}.
\label{eq:partition-sum}
\end{equation}
Moreover, for $1\le p\le k-1$ and admissible $h$,
\begin{align}
p(t_p=h\mid y,k)
&=\frac{L_{p,h}R_{k-p,h}}{L_{k,n}},\nonumber\\
p(b_h=1\mid y,k)
&=\sum_{p=1}^{k-1}p(t_p=h\mid y,k).
\label{eq:boundary-marginals}
\end{align}
\end{theorem}
\begin{proof}
\textbf{Step 1 (enumerate prefixes).}
For $k=0$ the unique empty path has weight one at $j=0$, and no empty path has a positive endpoint. Assume $L_{k,h}$ sums the local-score products over all $k$-segment partitions of $(0,h]$.
\textbf{Step 2 (split by the last segment).}
Every $(k+1)$-segment partition of $(0,j]$ has exactly one penultimate boundary $h\in\{k,\ldots,j-1\}$. Its weight is the weight of its $k$-segment prefix times $\widetilde A^{(0)}_{h,j}$. Summing over all such $h$ gives Equation~\eqref{eq:forward} and proves the prefix identity by induction.
\textbf{Step 3 (enumerate suffixes).}
The analogous decomposition by the first segment after $i$ gives Equation~\eqref{eq:backward}. At $i=0$, the suffix enumeration covers precisely the same set of full paths as the prefix enumeration at $j=n$. Hence $R_{k,0}=L_{k,n}$, proving Equation~\eqref{eq:partition-sum}.
\textbf{Step 4 (fix an ordered boundary).}
Under the event $t_p=h$, each admissible path is uniquely split into a $p$-segment prefix ending at $h$ and a $(k-p)$-segment suffix beginning there. Multiplication of their total sums gives the unnormalized numerator $L_{p,h}R_{k-p,h}$. The prior factor $C_k^{-1}$ is common to numerator and denominator and cancels, so division by $L_{k,n}$ proves the first boundary identity.
\textbf{Step 5 (forget the boundary rank).}
Within a strictly ordered partition, at most one of the events $\{t_p=h\}$ can occur at a given interior index $h$. They are disjoint as $p$ varies. Their probabilities therefore add, yielding the boundary-event marginal in Equation~\eqref{eq:boundary-marginals}. Normalization over $h$ of each ordered-boundary distribution gives one; summation over $h$ of boundary-event probabilities gives $k-1$.
\end{proof}

Let $p(k)$ be a segment-count prior. Because $C_k(u)$ normalizes Equation~\eqref{eq:partition-prior},
\begin{equation}
 p(k\mid y)\propto p(k)\frac{L_{k,n}}{C_k(u)},
 \qquad k=1,\ldots,k_{\max}.
\label{eq:k-posterior}
\end{equation}
The posterior may be reported in full; choosing $\widehat k=\arg\max_k p(k\mid y)$ is a separate decision rule. Conditional boundary distributions in Equation~\eqref{eq:boundary-marginals} and boundary-event probabilities can then be computed at fixed $k$ or averaged over $p(k\mid y)$.

\subsection{Posterior moments and Bayes curves}
For $r\in\{0,1,2\}$, define a block-covering contribution at fixed $k$,
\begin{equation}
F^{(r)}_{ij}(k)
=\frac{1}{L_{k,n}}
  \sum_{m=1}^{k}L_{m-1,i}\widetilde A^{(r)}_{ij}R_{k-m,j}.
\label{eq:block-cover}
\end{equation}
When the target at coordinate $s$ is the segment-level quantity $m_\star(\theta)$, its conditional posterior moments are
\begin{equation}
\mathbb E\!\left[m_\star(\theta_s)^r\mid y,k\right]
 =\sum_{i<s\le j}F^{(r)}_{ij}(k).
\label{eq:bayes-curve}
\end{equation}
The posterior mean Bayes curve uses $r=1$; the $r=2$ result gives pointwise posterior variance after subtracting the squared mean.

\begin{proposition}[Block-covering decomposition]
\label{prop:block-cover}
If every admissible partition assigns each coordinate to exactly one block and $A^{(r)}_{ij}$ is the integrated $r$th moment numerator for that block, then Equation~\eqref{eq:bayes-curve} is the exact fixed-$k$ posterior moment of the segment-level target at $s$.
\end{proposition}
\begin{proof}
\textbf{Step 1 (uniqueness of the covering block).}
Fix $s\in\{1,\ldots,n\}$. In every admissible partition there is exactly one segment $(i,j]$ for which $i<s\le j$; the corresponding segment position $m$ is also unique.
\textbf{Step 2 (integrate a fixed covering block).}
Conditioning on $(i,j]$ occupying position $m$, all earlier segments contribute the prefix sum $L_{m-1,i}$. The covering block contributes $\widetilde A^{(r)}_{ij}$, because its signal functional is raised to the $r$th power before integration. All later segments contribute $R_{k-m,j}$.
\textbf{Step 3 (sum without duplication).}
Summing over positions $m$ and blocks $(i,j]$ covering $s$ enumerates each full $k$-segment path once, multiplied by its integrated target contribution. This gives the numerator of the posterior expectation.
\textbf{Step 4 (normalize).}
The denominator $L_{k,n}$ is the sum of unweighted path contributions by Theorem~\ref{thm:sum-product}. Division proves Equation~\eqref{eq:bayes-curve}. When $r=2$, subtracting the square of the $r=1$ result yields pointwise posterior variance, with no claim of a simultaneous credible band.
\end{proof}

\subsection{Joint MAP segmentation}
Define local log scores $\ell_{ij}=\log\widetilde A^{(0)}_{ij}$ and max-sum messages
\begin{equation}
M_{k,j}=\max_{h\in\{k-1,\ldots,j-1\}}
       \{M_{k-1,h}+\ell_{h,j}\},
\quad M_{0,0}=0,
\label{eq:max-sum}
\end{equation}
with $M_{0,j}=-\infty$ for $j>0$. Store a deterministic maximizing predecessor $h^\star(k,j)$. Starting at $(k,n)$ and following these predecessors recovers a complete boundary vector.

\begin{theorem}[Correctness of max-sum backtracking]
\label{thm:map-correctness-paper}
For fixed $k$, the boundary vector returned by Equation~\eqref{eq:max-sum} and its stored predecessors belongs to
\begin{equation}
\arg\max_{0=t_0<\cdots<t_k=n}
  \sum_{q=1}^{k}\ell_{t_{q-1},t_q}.
\label{eq:joint-map-objective}
\end{equation}
Consequently, it is a joint MAP segmentation conditional on $k$. Ties are resolved only by the declared predecessor convention.
\end{theorem}
\begin{proof}
\textbf{Step 1 (base case).}
For a one-segment prefix $(0,j]$, the only admissible path consists of the block $(0,j]$. Hence $M_{1,j}=\ell_{0,j}$ is its maximum score, and its stored predecessor is $0$.
\textbf{Step 2 (induction hypothesis).}
Assume that for each feasible $h$, $M_{k-1,h}$ is the maximum log score among all $(k-1)$-segment partitions ending at $h$.
\textbf{Step 3 (optimal substructure).}
Every $k$-segment path ending at $j$ has a unique penultimate endpoint $h\in\{k-1,\ldots,j-1\}$. Its score is the prefix path score plus $\ell_{h,j}$. At any fixed $h$, the best possible prefix is $M_{k-1,h}$ by the hypothesis. Therefore maximizing over $h$ gives exactly Equation~\eqref{eq:max-sum}, the best $k$-segment score.
\textbf{Step 4 (recover a feasible maximizing path).}
Each predecessor $h^\star(k,j)$ attains the recurrence maximum and satisfies $h^\star(k,j)<j$. Starting from $(k,n)$ and following these predecessors decreases the endpoint and block count by one each step until $(0,0)$. The resulting ordered boundary vector attains $M_{k,n}$ and maximizes Equation~\eqref{eq:joint-map-objective}. This is conditional on $k$; joint optimization over $k$ also needs the count-prior offset.
\end{proof}

\begin{cautionbox}[Marginal modes are not a joint MAP path]
The vector formed from $\arg\max_h p(t_p=h\mid y,k)$ for each $p$ optimizes separate marginal decisions. Its entries can be incompatible or can have lower joint posterior probability than another partition. \BayesBreak therefore reports marginal modes as diagnostics and reserves ``joint MAP'' for Theorem~\ref{thm:map-correctness-paper}.
\end{cautionbox}

\subsection{Decision-theoretic segmentation summaries}
\label{sec:decision-rules}
Posterior summaries answer different decision questions. The distinction is not a formatting convention: the marginal posterior of $k$ sums over segmentations, while joint MAP scoring maximizes over them.

\begin{proposition}[Three non-equivalent posterior modes]
\label{prop:three-modes}
Define
\begin{align*}
\widehat k_{\rm marg}&\in\arg\max_k\left\{\log p(k)+\log L_{k,n}-\log C_k\right\},\\
\widehat t_{\rm cond}&\in\arg\max_{t\in\mathcal T_{\widehat k_{\rm marg}}}
 \sum_{q=1}^{\widehat k_{\rm marg}}\ell_{t_{q-1},t_q},\\
(\widehat k_{\rm joint},\widehat t_{\rm joint})&\in\arg\max_{k,t\in\mathcal T_k}
\left\{\log p(k)-\log C_k+\sum_{q=1}^{k}\ell_{t_{q-1},t_q}\right\}.
\end{align*}
Here $\ell_{ij}=\log\widetilde A^{(0)}_{ij}$. The first is a marginal-MAP segment count, the second a conditional joint-MAP partition at that count, and the third a global joint-MAP count--partition pair. In general they are distinct decisions.
\end{proposition}
\begin{proof}
\textbf{Step 1 (marginalize partitions).} The evidence for a fixed count is $p(y\mid k)=L_{k,n}/C_k$, so Bayes' rule gives the first maximization after discarding its constant normalizer.
\textbf{Step 2 (optimize within a selected count).} Conditional on $k$, the factors $p(k)$ and $C_k^{-1}$ are constant across admissible partitions. Therefore max-sum backtracking solves the second maximization for $k=\widehat k_{\rm marg}$.
\textbf{Step 3 (optimize jointly).} If $k$ also varies, the complete joint score retains both $\log p(k)$ and $-\log C_k$, giving the third maximization.
\textbf{Step 4 (show the distinction).} For each $k$, the marginal score uses a \emph{sum} $L_{k,n}$ over all paths, whereas the best-path score uses their \emph{maximum}. Counts with many moderately probable paths can have greater marginal support than counts with one sharply preferred path, so no identity between the maximizing counts follows.
\end{proof}

\begin{proposition}[Bayes actions for full-partition and coordinate losses]
\label{prop:loss-bayes-action}
For fixed $k$, 0--1 loss on the entire ordered vector $t=(t_1,\ldots,t_{k-1})$ is minimized by a joint MAP partition. For coordinatewise Hamming loss $\mathcal L(a,t)=\sum_{q=1}^{k-1}\mathbf1\{a_q\ne t_q\}$, constrained to admissible ordered actions $a_1<\cdots<a_{k-1}$, the Bayes action instead solves
\begin{equation}
\max_{1\le a_1<\cdots<a_{k-1}\le n-1}
\sum_{q=1}^{k-1}p(t_q=a_q\mid y,k).
\label{eq:constrained-hamming}
\end{equation}
This constrained optimization has its own additive dynamic program; independent coordinatewise marginal modes need not form an admissible action.
\end{proposition}
\begin{proof}
\textbf{Step 1 (0--1 loss).} Its posterior risk at action $a$ is $1-p(t=a\mid y,k)$, minimized exactly by maximizing the joint posterior mass.
\textbf{Step 2 (Hamming loss).} By linearity of conditional expectation, its posterior risk is $(k-1)-\sum_qp(t_q=a_q\mid y,k)$. Minimizing risk over valid actions is therefore equivalent to Equation~\eqref{eq:constrained-hamming}; dropping the ordering constraint would solve a different optimization problem.
\textbf{Step 3 (dynamic program).} Set $D_{1,h}=p(t_1=h\mid y,k)$ and recursively
\[
D_{q,h}=p(t_q=h\mid y,k)+\max_{g<h}D_{q-1,g}
\]
on valid index ranges, retaining predecessors. Backtracking from the best valid terminal state yields the constrained Bayes action; prefix maxima implement each layer efficiently.
\textbf{Step 4 (explicit counterexample).} For $n=7,k=3$, assign posterior masses $0.17,0.18,0.20,0.20,0.125,0.125$ to the admissible pairs $(1,3),(2,3),(4,5),(4,6),(1,5),(2,6)$, respectively, and zero elsewhere. They sum to one. The marginal mode of $t_1$ is $4$ with probability $0.40$, while that of $t_2$ is $3$ with probability $0.35$, producing the impossible vector $(4,3)$. The joint modes are $(4,5)$ and $(4,6)$, each of mass $0.20$; the constrained Hamming calculation must respect ordering rather than export $(4,3)$.
\end{proof}

The output interface reports each object by its complete name. Posterior-mean segment levels \emph{conditional on} an exported MAP partition are likewise distinct from posterior-mean Bayes curves averaged over the segmentation distribution.


\subsection{Complexity and numerical invariants}
If cumulative sufficient statistics make each local conjugate score constant time, there are $\Theta(n^2)$ blocks to precompute. Sum-product and max-sum each inspect $\Theta(k_{\max}n^2)$ predecessor transitions. Values and backpointers require $\Theta(k_{\max}n)$ working storage, while retaining all block-cover terms can require $\Theta(k_{\max}n^2)$ storage.

\begin{proposition}[Complexity]
\label{prop:complexity-paper}
When cumulative sufficient statistics make each conjugate segment calculation constant time, segment precomputation costs $\Theta(n^2)$ time and storage; either partition recursion costs $\Theta(k_{\max}n^2)$ time. For a nonconjugate block routine with per-block cost $T_{\rm block}$, local precomputation costs $\Theta(n^2T_{\rm block})$ and the global recursions are unchanged.
\end{proposition}
\begin{proof}
\textbf{Step 1 (candidate-segment count).}
There are $n(n+1)/2$ nonempty contiguous segments $(i,j]$ with $0\le i<j\le n$. With constant-time evaluation from prefix statistics, filling the triangular local score array therefore takes $\Theta(n^2)$ time and memory.
\textbf{Step 2 (count predecessor transitions).}
At a feasible state $(k,j)$, sum-product or max-sum examines $j-k+1$ predecessor endpoints in the unconstrained dense case. For a range of counts $1\le k\le k_{\max}\le n$, the exact transition total is of order $\sum_{k=1}^{k_{\max}}\sum_{j=k}^{n}(j-k+1)$, which is $\mathcal O(k_{\max}n^2)$ and attains that order for the usual regimes where $k_{\max}$ does not collapse to an empty or singleton state set.
\textbf{Step 3 (separate local approximation cost).}
If each local score costs $T_{\rm block}$ operations rather than a constant, filling the array costs $\Theta(n^2T_{\rm block})$. Neither sum-product nor max-sum changes its predecessor traversal as a result. These are worst-case accounting statements; finite-range measured timing slopes need not equal the asymptotic exponents.
\end{proof}

Four numerical identities are checked in every implementation: the normalized segment-count posterior sums to one; forward and backward total evidence agree; conditional boundary-event probabilities sum to $k-1$; and the backtracked path score equals the terminal max-sum score within numerical tolerance. These checks test algebraic implementation rather than empirical model quality.
